\documentclass[11pt]{article}
\usepackage{amsmath,amssymb,amsthm}
\newtheorem{theorem}{Theorem}
\newcommand{\E}{\mathbb E}
\newcommand{\Cov}{\operatorname{Cov}}
\newcommand{\Var}{\operatorname{Var}}
\title{Actual marginal Wilson scores and a quantitative inherited-mean interface}\author{}\date{}
\begin{document}\maketitle
\paragraph{Exact density after a coordinate elimination.}
Partition the original bonds into retained $R$ and eliminated $E$, and observe
$q(U)=U_R=y$. With the original normalized Haar factors the marginal density is
\[
 \rho_R(y)=Z^{-1}\int e^{-\beta S_0(y,z)}d\lambda_E(z),\qquad
 S_R^{\rm eff}(y)=-\log\rho_R(y).
\]
The positive denominator is the actual marginal integral. No determinant or coordinate
Jacobian is introduced by this coordinate projection. If charts are later used, their actual
Haar Jacobians remain part of that density. All conditional means below use this exact integral.
For a deterministic invariant direction $X_v$ on retained links, write $F_v=X_vS_0$.
Differentiation under the compact integral gives
\begin{align}
 X_vS_R^{\rm eff}&=\beta\E(F_v\mid y),\label{score}\\
 X_v^2S_R^{\rm eff}&=\beta\E(X_v^2S_0\mid y)-\beta^2\Var(F_v\mid y),\label{hessian}\\
 X_v\E(G\mid y)&=\E(X_vG\mid y)-\beta\Cov(G,F_v\mid y).\label{response}
\end{align}
In particular the effective score is not the remaining bare Wilson score.
These identities retain the response of the genuine denominator as the centered covariance.

\paragraph{Dimension-independent directional Ward moments.}
Let $v$ have arbitrary finite support and constant Lie coefficients, with
$\|v\|_2^2=\sum_{b,A}|v_b^A|^2$. The compact plaquette action has uniformly bounded
second derivatives. Bounded plaquette incidence therefore implies
\[
 |X_v^2S_0|\le C_N\|v\|_2^2
\]
pointwise, independently of the number of bonds in the support.
Indeed each plaquette contributes at most $C(\sum_{b\in p}|v_b|)^2$ and each bond
belongs to exactly six plaquettes. Since $X_v$ is Haar divergence-free, the same recurrence
as the uniform contact moments calculation proves
\[
 \E|F_v|^{2m}\le(2m-1)!!(C_N\|v\|_2^2/\beta)^m.
\]
Thus for every fixed $p\ge1$, under the genuine marginal $\nu_R$,
\begin{align}
 \|X_vS_R^{\rm eff}\|_{L^p(\nu_R)}&\le C_{N,p}\sqrt\beta\|v\|_2,\label{scorebound}\\
 \|X_v^2S_R^{\rm eff}\|_{L^p(\nu_R)}&\le C_{N,p}\beta\|v\|_2^2,\label{hessbound}\\
 X_v^2S_R^{\rm eff}&\le C_N\beta\|v\|_2^2\quad\hbox{pointwise}.\label{onesided}
\end{align}
For \eqref{hessbound}, use conditional Jensen on $\Var(F_v|y)\le\E(F_v^2|y)$;
for \eqref{onesided}, keep the negative variance sign. Polarization gives the analogous
$C_p\beta\|v\|_2\|w\|_2$ bound for the symmetrized mixed Hessian.
An ordered mixed Hessian has in addition half the commutator score, which is bounded
by \eqref{scorebound} and $\|[v,w]\|_2\le C\|v\|_2\|w\|_2$.
These estimates hold for ANY eliminated coordinate set. They are averaged directional
estimates, not a pointwise lower Hessian bound, a spectral gap, or a summable kernel bound.

\paragraph{Proved finite range for the first separated-link layer.}
For the noninteracting free links of AI36, no plaquette contains two eliminated links.
Write $S_0=S_{\rm out}(y)+\sum_{e\in E}s_e(U_e;y)$ with every touching plaquette
assigned to its unique free link, and put
$z_e(y)=\int e^{-\beta s_e(U_e;y)}dU_e$. Then exactly
\[
 \rho_R(y)=Z^{-1}e^{-\beta S_{\rm out}(y)}\prod_{e\in E}z_e(y),\qquad
 S_R^{\rm eff}=\log Z+\beta S_{\rm out}-\sum_{e\in E}\log z_e.
\]
Each $z_e$ depends only on the literal outside plaquette star of $e$.
Consequently every score and Hessian entry has a fixed finite-range carrier, and the
number of such carriers meeting a retained bond is bounded by lattice incidence.
This is actual locality of the effective action for this layer. For a general eliminated
set the covariance in \eqref{hessian} may connect far-separated retained bonds; the
moment estimates above do not claim it vanishes or decays.

\paragraph{A derivative bound for the actual inherited H mean.}
Take the separated packing and $G(f)=O_H(f)$ or the specified $r_H(f)$.
Let $\overline G(f)=\E(G(f)|y)$. For a retained bond/color derivative $X_b^A$, let
$\mathcal N_b$ consist of the finitely many original force rows whose terms either depend
directly on $U_b$ or vary in an eliminated block whose plaquette star meets $b$.
Set $W_b(f)^2=\sum_{c\in\mathcal N_b}|f_c|^2$, retaining actual pointed offsets when
present. The number of these rows and their multiplicity over $b$ are fixed constants.
Then
\begin{equation}\label{meanbound}
 \|X_b^A\overline G(f)\|_{L^2(\nu_R)}
 \le C\sigma\beta^{-1/2}\log^{3/2}(2+\beta)\,W_b(f),
 \qquad \sigma=g^{-2}a^{-2}.
\end{equation}
Here is a proof that uses the actual conditional covariance rather than an unrelated norm.
The direct derivative $X_b^AG$ is a fixed sum of force times bounded force derivative,
plus the literal Casimir derivative for $r_H$. The plaquette energy bound gives
$\|X_b^AG\|_2\le C\sigma\sqrt{L_\beta/\beta}\,W_b(f)$, with
$L_\beta=\log(2+\beta)$.

Conditional product independence localizes $\Cov(G,F_b^A|y)$ to those blocks sharing
free variables with $F_b^A$. Write $G_b$ for exactly their varying source terms and
$v_b=\Var(G_b|y)$, $c_b=\Var(F_b^A|y)$.
The local ordinary covariance calculation proves $\E v_b\le C\sigma^2\beta^{-2}L_\beta^2W_b(f)^2$.
Compactness of this fixed local sum also gives $v_b\le C\sigma^2W_b(f)^2$ pointwise.
The directional Ward moments give $\|c_b\|_m\le Cm/\beta$ for integers $m\ge2$.
Conditional Cauchy--Schwarz and Holder thus yield
\[
 \|\Cov(G_b,F_b^A|y)\|_2^2
 \le\E(v_bc_b)
 \le\frac{Cm}{\beta}
       (C\sigma^2W_b^2)^{1/m}
       (C\sigma^2\beta^{-2}L_\beta^2W_b^2)^{1-1/m}.
\]
Choosing $m=\max(2,\lceil L_\beta\rceil)$ bounds this by
$C\sigma^2\beta^{-3}L_\beta^3W_b^2$.
Insert this and the direct derivative estimate in \eqref{response} to obtain
\eqref{meanbound}. A bounded physical boundary variation $U_b\mapsto e^{aZ_b}U_b$
adds its original factor $a$ to this derivative. Saved field-independent source factors
multiply the whole expression; field-dependent ones remain inside $G$ and are differentiated.

Summing actual derivative components gives the quantitative finite-range estimate
\[
 \E_{\nu_R}\sum_{b,A}|X_b^A\overline G(f)|^2
 \le C\sigma^2\beta^{-1}L_\beta^3\sum_c|f_c|^2.
\]
The final sum can be restricted to the literal completed support. This is a Dirichlet-energy
bound for the inherited conditional mean, not a variance bound without an additional
coercive or spatial covariance input.

\paragraph{Selectors, current partners, and denominator scope.}
For a bath-measurable selector $\chi(y)$, the selected marginal is $\chi\rho_R$.
A smooth positive selector adds $-X_v\log\chi$ to the effective score.
A sharp selector must instead retain the weak derivative
$X_v(\chi\rho_R)=(X_v\chi)\rho_R+\chi X_v\rho_R$; there is no pointwise log score
on its boundary. For a selector depending on eliminated variables, the undivided exact
formula is
\[
 X_v\int\chi(y,z)e^{-\beta S_0(y,z)}dz
 =\int(X_v\chi-\beta\chi F_v)e^{-\beta S_0}dz.
\]
Its normalized form requires its genuine nonzero selected denominator and its actual
conditional measure. No preceding probability contraction supplies a bound for that extra
response. The same rule applies to a live kernel, constraint, or coordinate Jacobian.
The AI41 current supported on eliminated bonds still has $Dq\mathcal W=0$ for this
observation, as proved in the actual AI36 interface calculation. Outer currents act on the complete marginal and hence
on every $z_e$; their responses are exactly the effective-score terms above. No boundary
partner or independent absolute scalar factor is canceled by these estimates.

\paragraph{Next interface.}
We now have actual marginal score moments, first-layer finite-range effective action, and
an $L^2$ derivative/locality bound for the inherited ordinary mean. To conclude a bound
on $\Var_{\nu_R}\overline G$ from its Dirichlet energy would require a uniform Poincare
or appropriate restricted covariance estimate for this actual marginal. Such an estimate
has not been proved. The one-sided Hessian upper bound \eqref{onesided} does not provide
it. Subsequent coordinate marginals retain \eqref{scorebound}--\eqref{hessbound}, but
can lose the first-layer finite range and can generate long-range conditional-mean covariance.
No general MT36 or universal operator normalization is claimed.
\end{document}
